% =========================================================
% Partial and Total Maps
% =========================================================

\section{Partial and Total Maps}
\label{sec:partial-total-maps}

\begin{remark*}[Toolkit: Partial and Total Maps]
Before studying maps that preserve order structure, the foundational
function types are recorded: a partial map need not be defined everywhere,
while a total map is defined on every input.  These are the base layer
upon which the order-theoretic notions rest.
\end{remark*}

% ---------------------------------------------------------
% def:partial-map
% ---------------------------------------------------------

\begin{definitionbox}{Definition (Partial Map)}
\begin{definition}[Partial Map]
\label{def:partial-map}
Let $A$ and $B$ be sets.  A \emph{partial map} $f : A \rightharpoonup B$
is a relation $f \subseteq A \times B$ such that every element of $A$ has
at most one image:
\[
\forall x \in A,\; \forall y_1, y_2 \in B,\quad
(x, y_1) \in f \;\land\; (x, y_2) \in f
\;\Longrightarrow\; y_1 = y_2.
\]
\end{definition}
\end{definitionbox}

\begin{remark*}[Standard quantified statement]
\[
\operatorname{PartialMap}(f, A, B)
\;\coloneqq\;
f \subseteq A \times B
\;\land\;
\forall x \in A,\; \forall y_1, y_2 \in B,\;
(x,y_1) \in f \land (x,y_2) \in f \Rightarrow y_1 = y_2.
\]
\end{remark*}

\begin{remark*}[Interpretation]
A partial map assigns to each element of $A$ at most one element of $B$.
Its domain of definition
$\operatorname{dom}(f) = \{x \in A : \exists y \in B,\, (x,y) \in f\}$
may be a proper subset of $A$.  Partial maps arise in order theory when
a function is well-defined only on elements satisfying some order-theoretic
condition --- for example, on elements that admit an upper bound.
\end{remark*}

\NoLocalDependencies
% ---------------------------------------------------------
% def:total-map
% ---------------------------------------------------------

\begin{definitionbox}{Definition}
\begin{definition}[Total Map]
\label{def:total-map}
A \hyperref[def:partial-map]{partial map} $f : A \rightharpoonup B$ is a
\emph{total map} if it is defined on every element of $A$:
\[
\forall x \in A,\; \exists y \in B,\quad (x, y) \in f.
\]
A total map is written $f : A \to B$ and is precisely a function in the
usual sense.
\end{definition}
\end{definitionbox}

\begin{remark*}[Standard quantified statement]
\[
\operatorname{TotalMap}(f, A, B)
\;\coloneqq\;
\operatorname{PartialMap}(f, A, B)
\;\land\;
\forall x \in A,\; \exists y \in B,\; (x, y) \in f.
\]
\end{remark*}

\begin{dependencies}
\hyperref[def:partial-map]{Partial map},
\hyperref[def:function]{Function}.
\end{dependencies}
